\section{Introduction}

Modern LLM-powered coding assistants --- Claude Code, GitHub Copilot, Cursor,
Gemini CLI, and others --- answer developer questions by issuing sequences of
file reads and search calls, accumulating raw source text until the question can
be answered. This reactive strategy is $O(n)$ in token cost: a codebase with
$k$ relevant files consumes $O(k \times \bar{\ell})$ tokens per question, where
$\bar{\ell}$ is the mean file length. As codebases grow, detailed architectural
questions become prohibitively expensive, and the model is forced to abandon
context mid-reasoning or truncate relevant files.

Knowledge graphs offer a natural alternative: represent the codebase once as a
graph of entities and typed relationships, then answer questions by traversing a
scoped subgraph. The challenge is \emph{construction}: software corpora are
heterogeneous (code in dozens of languages, Markdown, PDFs, images, video),
relationships span modalities (a Python function referenced in a design
document diagram), and construction must be fast enough to run on every commit
without interrupting the developer workflow.

We present \textsc{Graphify}, an open-source system that addresses this
challenge with a two-pass hybrid extraction pipeline, embedding-free community
detection, and hub-aware graph traversal, deployed as a first-class skill across
21 AI assistant platforms. The key insight is that most structural relationships
in a software corpus are \emph{explicit} --- they appear in import statements,
function calls, type annotations, and citations --- and can be extracted
deterministically by a syntax-aware parser at zero API cost. Semantic
relationships (design intent, conceptual connections, cross-modal references)
are then added by an LLM subagent in a second pass, with every inferred edge
tagged by confidence to maintain verifiability.

Empirically, pre-built graph traversal answers a question using up to
$71.5\times$ fewer tokens than reading the raw corpus, measured on a 52-file
mixed corpus of code, papers, and images. On the same corpus, parallel AST
extraction achieves $3.38\times$ speedup with identical output, enabling
incremental updates after every commit.

\smallskip
\noindent\textbf{Contributions.} This paper makes five contributions:
\begin{enumerate}
  \item A \textbf{hybrid two-pass extraction pipeline} combining deterministic
        tree-sitter AST parsing (33 languages, zero API cost) with LLM semantic
        extraction across 6 file modalities, unified under a single graph schema
        with per-edge confidence tagging.
  \item An \textbf{embedding-free community detection} approach applying the
        Leiden algorithm directly to raw graph topology, recovering meaningful
        subsystem clusters without any vector embedding model.
  \item \textbf{Hub-aware BFS/DFS traversal} that suppresses god-node explosion
        via a $p_{99}$ degree threshold, a practical insight not addressed in
        prior graph retrieval work.
  \item \textbf{Quantitative evaluation} demonstrating up to $71.5\times$ token
        reduction, $89.9\%$ extraction precision, and $3.38\times$ parallel
        throughput across five real-world corpora.
  \item \textbf{Real-world validation at scale}: 1.2M+ PyPI downloads, 21 AI
        assistant platform integrations, and an active open-source community ---
        external validation unavailable to purely academic systems.
\end{enumerate}

The remainder of this paper is organized as follows.
Section~\ref{sec:related} surveys related work.
Section~\ref{sec:system} describes the Graphify architecture.
Section~\ref{sec:eval} reports experimental results.
Section~\ref{sec:conclusion} concludes.
